\documentclass[margin,line,pifont,palatino,courier]{res}
\usepackage{xcolor}
\usepackage{pifont}
\usepackage[utf8]{inputenc}
\usepackage{hyperref}
\usepackage{cite}
\usepackage{graphicx}
\textheight=9.0in
\usepackage{fancyhdr}
\usepackage{bibentry}
\pagestyle{fancy}
\renewcommand{\headrulewidth}{0pt}
\fancyhf{}
\rfoot{{\footnotesize Curriculum Vitae, Thejus Mahajan, \thepage}}

\newenvironment{list1}{
  \begin{list}{\ding{113}}{%
      \setlength{\itemsep}{0in}
      \setlength{\parsep}{0in} \setlength{\parskip}{0in}
      \setlength{\topsep}{0in} \setlength{\partopsep}{0in}
      \setlength{\leftmargin}{0.1in}
      }}{\end{list}}
\newenvironment{list2}{
  \begin{list}{$\bullet$}{%
      \setlength{\itemsep}{0in}
      \setlength{\parsep}{0in} \setlength{\parskip}{0in}
      \setlength{\topsep}{0in} \setlength{\partopsep}{0in}
      \setlength{\leftmargin}{0.1in}}}{\end{list}}

\begin{document}

\name{\hphantom{sss} Dr. Thejus Mahajan \vspace*{.1in}}

\begin{resume}

\section{\sc Contact Information}

\vspace{.05in}
\begin{tabular}{@{}p{3.75in}p{2in}}
 Email: \verb+thejus.mahajan@uni-hamburg.de+\\
 Ph: +49-15259697629 \\
 Website: \url{https://thejusmahajan.github.io} \\
 Hamburg, Germany
\end{tabular}

\section{\sc Career Interests}
To apply expertise in computational modeling, data science, and bioinformatics to address complex challenges in clinical and biological research. Skilled in multivariate statistical modeling, biostatistics, and developing analysis pipelines using R (Bioconductor, Tidymodels) and Python. Actively transitioning into clinical bioinformatics through specialized training on patient datasets and biostatistical methods for health research.

\section{\sc Work Experience}

\begin{itemize}
\item {\bf CQ Beratung + Bildung, Berlin, Germany}\\
\vspace*{-.2in}
\item[] Further Training -- Application-related Bioinformatics and Biostatistics\\
August 2025 -- Present
\vspace*{.1in}
\item[-] Shell (BASH), Python (object-oriented), and programming methodology.
\item[-] Bioinformatics databases (NCBI, Biopython, SQL).
\item[-] Sequence analysis, structural bioinformatics, and NGS analysis (Nextflow, Galaxy).
\item[-] Applied biostatistics with real and simulated patient datasets (R/Bioconductor, ANOVA, PCA, HCA, survival analysis).
\end{itemize}
\vspace*{.1in}

\begin{itemize}
\item {\bf HealthTwiSt GmbH, Berlin, Germany}\\
\vspace*{-.2in}
\item[] Internship -- Machine Learning in R\\
February 2026 -- April 2026
\vspace*{.1in}
\item[-] Implementation of Machine Learning workflows using the Tidymodels framework in R.
\item[-] Supervised by Dr. Andreas Busjahn (HealthTwiSt GmbH).
\item[-] Practical application of bioinformatics and biostatistical analysis skills.
\end{itemize}
\vspace*{.1in}

\begin{itemize}
\item {\bf Helmholtz-Zentrum Hereon, Geesthacht, Germany}\\
\vspace*{-.2in}
\item[] Guest Scientist -- Computational Biologist \& Systems Modeler\\
May 2025 -- October 2025
\vspace*{.1in}
\item[-] Spearheaded research modeling evolutionary processes in marine environments.
\item[-] Simulated long-term impacts of environmental changes on ecosystems.
\end{itemize}
\vspace*{.1in}

\begin{itemize}
\item {\bf University of Hamburg, Institute of Marine Ecosystem and Fisheries Sciences}\\
\vspace*{-.2in}
\item[] Post-doctoral Scientist -- Marine Ecosystem Modelling\\
August 2021 -- January 2025
\vspace*{.1in}
\item[-] Developed the ``Cyanobacteria Life Cycle'' (CLC) model within the ERGOM framework.
\item[-] Engineered models in Fortran and analyzed large NetCDF datasets using Python in Linux.
\item[-] Statistical analysis and visualization of ecological time-series data using R.
\end{itemize}
\vspace*{.1in}

\begin{itemize}
\item {\bf Universit\'e Paris-Saclay, Institut des Sciences Mol\'eculaires d'Orsay, France}\\
\vspace*{-.2in}
\item[] Ph.D. Researcher in Astrochemistry\\
October 2015 -- September 2018
\vspace*{.1in}
\item[-] Led atomic collision experiments and processed terabytes of raw data.
\item[-] Utilized C++ for signal processing and algorithm development.
\item[-] Developed a molecular fragmentation model for the KIDA database.
\item[-] {\bf Thesis}: ``Excitation and fragmentation of C${}_n$N${}^+$ (n=1-3) molecules in collisions with He atoms at intermediate velocity; fundamental aspects and application to astrochemistry.''
\item[-] {\bf Advisor}: Dr. Karine B\'eroff
\end{itemize}

\section{\sc Education}

\begin{itemize}
\item {\bf Ph.D. in Astrochemistry} (2015 -- 2018)\\
Institut des Sciences Mol\'eculaires (ISMO), Universit\'e Paris-Saclay, France
\vspace*{.1in}
\item {\bf M.Sc. in Physics} (July 2012 -- January 2015)\\
National Institute of Technology Calicut, India\\
CGPA: 8.2/10
\end{itemize}

\section{\sc Skills \& Competencies}

\begin{tabular}{@{}p{1.5in}p{4in}}
{\bf Programming:} & Python (Pandas, NumPy, Scikit-learn), R (Bioconductor, Tidymodels, ggplot2), Fortran, SQL, BASH\\[1ex]
{\bf Bioinformatics:} & NGS Analysis, Galaxy, Sequence Alignment, Structural Bioinformatics, Nextflow\\[1ex]
{\bf Biostatistics:} & Multivariate Analysis, PCA, ANOVA, Dimensionality Reduction, Markov Chains\\[1ex]
{\bf Data Analysis:} & Large dataset processing (NetCDF, HDF5), Data visualization, Pipeline development\\[1ex]
{\bf Tools:} & Linux/Unix, Git, HPC clusters, \LaTeX, Microsoft Office
\end{tabular}

\section{\sc Languages}

\begin{tabular}{@{}p{1.2in}p{4in}}
English: & Proficient (speaking, reading, writing)\\
German: & B1 (Goethe Certified)\\
French: & Intermediate\\
Malayalam: & Native\\
Tamil, Hindi: & Proficient\\
\end{tabular}

\section{\sc Training \& Courses}

\begin{itemize}
\item[-] HPC Training at J\"ulich Supercomputing Center (JSC)
\item[-] MITx: 7.00x Introduction to Biology -- The Secret of Life
\end{itemize}

\section{\sc Publications}

Google Scholar: \url{https://scholar.google.com/citations?hl=en&user=PJkZwAwAAAAJ}

\begin{enumerate}
\item T. Mahajan, et al. \textit{Journal of Physics: Conference Series}, 1412:142026, 2020.
\item T. Idbarkach, et al. \textit{Astronomy \& Astrophysics}, 628:A75, 2019.
\item T. Mahajan, et al. \textit{Journal of Physics B}, 2019.
\item T. Idbarkach, et al. \textit{Journal of Physics B}, 51(24):245201, 2018.
\item T. Idbarkach, et al. \textit{Journal of Physics: Conference Series}, 1412(11):112028, 2020.
\end{enumerate}

\section{\sc References}

\begin{itemize}
\item {\bf Dr. Andreas Busjahn}\\
HealthTwiSt GmbH, Berlin\\
Email: \texttt{andreas@busjahn.net}\\
Professional relation: Biostatistics instructor and internship supervisor

\item {\bf Prof. Dr. Kai Wirtz}\\
Helmholtz-Zentrum Hereon, \"Okosystemmodellierung\\
Email: \texttt{kai.wirtz@hereon.de}\\
Professional relation: Supervisor during Guest Scientist position

\item {\bf Prof. Dr. Elisa Schaum}\\
University of Hamburg, Institute of Marine Ecosystem and Fisheries Sciences\\
Email: \texttt{elisa.schaum@uni-hamburg.de}\\
Professional relation: Collaborator during Post-doctoral research
\end{itemize}

\end{resume}
\end{document}
